% Source PDF: source/普通高考/2013/2013重庆文.pdf, page 1; vector program flowchart.
% Grid evidence: tmp/review_2013_chongqing_liberal/grid/page-1-grid100.jpg and
%   q5_original_clean4.png (no-grid crop from the source page).
% Elements: rounded 开始/结束 terminals; initialization rectangle k=1,s=1;
% update rectangle s=s+(k-1)^2; increment rectangle k=k+1; condition diamond
% s>15?; output parallelogram 输出 k; solid arrows and 是/否 labels.
% Complete node -> directed-edge list from the source:
%   start -> init -> sum -> test; test --是--> output -> finish;
%   test --否--> increment -> return to the connector above sum. The return
%   arrow, main downward arrows, and output branch are independent; no other
%   edges or dashed segments are present. Node text uses local zero indentation,
%   centered alignment, and compact padding.

\begin{tikzpicture}[
  >=stealth,
  x=1cm,
  y=1.3cm,
  line width=0.55pt,
  line cap=round,
  line join=round,
  font=\normalsize,
  flowtext/.style={anchor=center,align=center,inner sep=0pt,
    execute at begin node={\setlength{\parindent}{0pt}}},
  every node/.style={flowtext},
  flowbox/.style={flowtext,draw,minimum height=0.68cm,inner xsep=4pt,inner ysep=2pt},
  terminal/.style={flowbox,rounded corners=0.20cm,minimum width=1.24cm},
  io/.style={flowbox,shape=trapezium,trapezium left angle=72,trapezium right angle=108,
    minimum height=0.76cm},
  decision/.style={flowbox,shape=diamond,aspect=2.8,inner xsep=4pt,inner ysep=2pt,
    minimum height=1.00cm}
]
  % Main vertical path.
  \node[terminal] (start) at (0,9.20) {开始};
  \node[flowbox,minimum width=3.55cm] (init) at (0,7.70) {$k=1,s=1$};
  \node[flowbox,minimum width=4.15cm] (sum) at (0,6.25) {$s=s+(k-1)^2$};
  \node[decision,minimum width=3.70cm] (test) at (0,4.78) {$s>15?$};
  \node[io,minimum width=2.55cm] (output) at (0,3.28) {输出 $k$};
  \node[terminal] (finish) at (0,1.84) {结束};
  \node[flowbox,minimum width=2.70cm] (inc) at (4.00,6.25) {$k=k+1$};

  \draw[->] (start.south) -- (init.north);
  \draw[->] (init.south) -- (sum.north);
  \draw[->] (sum.south) -- (test.north);
  \draw[->] (output.south) -- (finish.north);

  % The no branch goes right, up through k=k+1, then returns to the connector before the sum update.
  \draw[->] (test.east) -- node[flowtext,above] {否} (4.00,4.78) -- (inc.south);
  \draw[->] (inc.north) -- (4.00,7.00) -- (0,7.00);

  % The 是 label is attached to the terminating edge, not manually offset.
  \draw[->] (test.south) -- node[flowtext,right] {是} (output.north);
\end{tikzpicture}
